\newvideofile{Kirchhoff}{Die Kirchhoffschen Regeln}

\section{Kirchhoff's laws}




\begin{frame}{}
	\s{
	
		The component equations introduced in the previous module, which show the relationship between voltage 
        and current at the individual components, are not sufficient to calculate all voltages and currents 
        within a network. Kirchhoff's rules, also known as mesh or node rules, 
        provide the equations required for this.
	}
	\ftx{Learning objectives: Kirchhoff's rules}
	
	\begin{Learning objectives}{Kirchhoff's rules}
		The students can
		\begin{itemize}
			\item reproduce the key statements of Kirchhoff's rules
			\item apply Kirchhoff's rules to simple resistance networks

		\end{itemize}
	\end{Learning objectives}
	
	\speech{folie9}{1}{Im vorangegangenen Video wurde der Grundstromkreis mit seinen Maschen, Knoten und Zweigen eingeführt.
		Jetzt stellt sich die Frage, wie genau sich Strom und Spannung im Netzwerk aufteilen.
		Dafür reichen die Bauteilgleichungen allein nicht aus. 
		Genau hier kommen die Kirchhoffschen Regeln ins Spiel.
		Diese liefern uns die nötigen Gleichungen, um alle Spannungen und Ströme im Netzwerk berechnen zu können.
		Man unterscheidet dabei zwei zentrale Gesetze, nämlich die Knotenregel zur Beschreibung von Stromverzweigungen, und die Maschenregel zur Beschreibung von Spannungen
		in geschlossenen Pfaden.
		Nach dem Durcharbeiten dieses Kapitels solltet ihr in der Lage sein, 
		die Kernaussagen der Kirchhoffschen Regeln wiederzugeben, sowie
		sie auf einfache Widerstandsnetzwerke anzuwenden.}
		\silence{4}
	
\end{frame}



\begin{frame}
	\ftx{Node rule I (Kirchhoff's 1st law)}
	
	\s{
	
	
	
		\subsection{Node rule I (1. Kirchhoff rule)}
		
        The node rule states that the sum of all currents flowing into a node is identical to the sum of all currents flowing out of it:		
		
		\begin{Key point}{}
			
			\begin{center}
				Sum of incoming currents\\
				=\\
				Sum of outgoing currents\\
			\end{center}
			
		\end{Key point}
		
		Mathematically expressed, this results in the following relationship:
		\begin{equation}
			\sum I_\mathrm{zu} = \sum I_\mathrm{ab}
		\end{equation}
		
		It should be noted that the currents are evaluated according to their counting arrow direction.
		Currents flowing into the node are counted as positive, currents flowing out of the node
		are counted as negative. If the actual direction of a current is not known in advance,
		the counting arrow direction can be set arbitrarily. The actual current direction is determined
		from the numerical value as the result of the calculation. If this is negative for a current, i.e. $I < 0$, 
		the real current flows in the opposite direction.\\

		Alternatively, this relationship can also be expressed by the fact that the sum of all currents in a node is equal to 0:
		
		
		\begin{equation}
			\sum_{i=1}^{n} = 0
		\end{equation}
		
        Applied to the example nodes in Figure \ref{fig:knoten}, this results in the following equations:
		
		
		\begin{figure}[h!]
			\renewcommand{\figurename}{Figure}
			\centering
			{\includesvg[width=0.25\textwidth]{knotenregel1}}
			
			{\caption{Simple example node}}
			\label{fig:knoten}
		\end{figure}
		
		
		
		
		
		\begin{equation*}
			I_1 = I_2 + I_3
		\end{equation*}
		\begin{equation*}
			I_1 - I_2 - I_3 = 0
		\end{equation*}
		
		However, the node to which the rule refers does not have to consist of just one point.
		Instead, it is possible to define so-called \textbf{envelope nodes}, which completely enclose an area
		within a circuit.
		Applied to the envelope node shown in Figure 3.2
		%\ref{fig:knoten2},
		this results in the equation
		
		\begin{equation*}
			I_1 + I_2 - I_3 - I_4 = 0
		\end{equation*}
		
		
		\begin{figure}[h!]
			\renewcommand{\figurename}{Figure}
			\begin{center}
				\input{Tikz/src/2_kirchhoff/Huellknoten_Skript}
			\end{center}
			\label{fig:knoten2}
			\caption{Terminal node of an electrical circuit}
		\end{figure}
		
		The node rule is not only applicable to discrete components due to the generalisation of the charge conservation law for source-free
		flow fields. Rather, it can be applied to any real
		structure. In Figure \ref{fig:huell}, the current flowing into the wire
        $I$ is therefore equal to the total current flowing out of the envelope surface of the sheet metal cut-out.
		
		
		
		\begin{figure}[h!]
			\renewcommand{\figurename}{Figure}
			\centering
			\s{\includesvg[width=0.6\textwidth]{blechausschnitt}}
			
			\s{\caption{Wire connected to a sheet metal cut-out. The current $I$ flows into the wire, while a
				current density $\vec{J}$ flows out of the sheet metal cut-out.}}
			\label{fig:huell}
		\end{figure}
		
		In general, this relationship can be described by the closed surface integral over the envelope surface $\vec{A}$:
		
		\begin{equation}
			\iint_A \vec{J} \cdot d\vec{A} = 0 
		\end{equation}
		
		
		%\begin{equation*}
		%	I_1 = I_2 + I_3
		
		
		%\end{equation*}
		
		
		
		
	}
	
	\b{
	
		\begin{columns}
			\column[t]{0.58\textwidth}
			
			\vspace{-20pt}
			\begin{Key point}{}
				\begin{center}
					Total inflows\\
					
					=\\
					
					Total outflows\\
				\end{center}
				
			\end{Key point}
			
			%	\vspace{5pt}
			
			
			\begin{equation*}
				\sum I_\mathrm{in} = \sum I_\mathrm{out}
			\end{equation*}
			
			\phantom{.}\\
			
			
			
			
			\only<2->{
			
				Alternatively: count incoming currents as positive and outgoing currents as negative:\\
				
				
				\begin{equation*}
					\sum_{k=1}^{n}I_k = 0
				\end{equation*}
			}
			\column[t]{0.42\textwidth}
			\begin{figure}[h!]
				\centering
				\includesvg[width=0.7\textwidth]{knotenregel1}
			\end{figure}
			
			
			$I_1 = I_2 + I_3$\\
			
			\phantom{text}\\
			\pause
			$I_1-I_2-I_3=0$
		\end{columns}
		\speech{folie10}{1}{Starten wir mit der ersten Kirchhoffschen Regel, der Knotenregel.
			Sie besagt,
			Der gesamte elektrische Strom, der in einen Knoten reinfließt, muss auch wieder rausfließen.
			Oder mathematisch ausgedrückt, 
			Die Summe der zufließenden Ströme ist gleich der Summe der abfließenden Ströme.
			Verdeutlicht wird dies im nebenstehenden Beispiel. Der Strom I eins fließt in den Knoten hinein und befindet sich dementsprechend auf der linken Seite der Gleichung, 
			während die Ströme I zwei und I drei aus dem Knoten herausfließen und somit auf der rechten Seite der Gleichung stehen.}
		\silence{2}
		\speech{folie10}{2}{ Alternativ kann man diesen Zusammenhang auch so ausdrücken,
			Die Summe aller Ströme am Knoten ist null, wenn man das Vorzeichen der Ströme beachtet.
			Dabei gilt,
			In den Knoten hineingehende Ströme werden positiv gezählt,
			aus dem Knoten herausfließende Ströme hingegen negativ.
			Sollte man die tatsächliche Richtung des Stromes nicht kennen,so setzt man den Zählpfeil willkürlich.
			Falls man dieses Vorzeichen entgegen der tatsächlichen Stromflussrichtung gesetzt hat, resultiert dies rechnerisch in einem negativen Wert des Stromes, was für uns allerdings auch kein Problem darstellt.}
	}
	\silence{4}
	
\end{frame}

\begin{frame}
	\b{
		\ftx{Knot rule II}
		
		\begin{columns}
			\column[t]{0.5\textwidth}
			Application to envelope nodes:
			
			\vspace{-5pt}
			
			\begin{figure}[h!]
				\begin{center}
					\input{Tikz/src/2_kirchhoff/Huellknoten_Folien}
				\end{center}
				
			\end{figure}
			\vspace{-12pt}
			
			\begin{equation*}
				I_1 + I_2 - I_3 - I_4 = 0
			\end{equation*}
			
			
			
			\column[t]{0.5\textwidth}
			\pause
			
			Application to geometric structures:
			
			\begin{figure}[h!]
				\centering
				\includesvg[width=0.9\textwidth]{blechausschnitt}
			\end{figure}
			
			\begin{equation*}
				\iint_A \vec{J} \cdot d\vec{A} = 0 
			\end{equation*}
			
			
			
		\end{columns}
		\speech{folie11}{1}{
			Ein solcher Knoten muss nicht nur ein einzelner Punkt sein. Auch ein Bereich der Schaltung, welcher von einer Linie vollständig umschlossen wird, kann als ein sogenannter 
			Hüllknoten betrachtet werden. Auch hier gilt die Regel, dass der Strom, der irgendwo in diesen Hüllknoten hereinfließt, irgendwo wieder herausfließen muss.
			In diesem Beispiel stellt der grüne Kreis den Hüllknoten dar. Dieser enthält eine Schaltung aus vier Widerständen. In diese Schaltung fließen die Ströme I eins sowie I zwei.
			Hingegen verlassen die Ströme I drei und I vier den Bereich. 
			Folglich werden die Ströme eins und Zwei mit positivem Vorzeichen gezählt, während I drei und I vier ein negatives Vorzeichen haben. Die Summe dieser Ströme ist null.}
			\silence{2}
		\speech{folie11}{2}{
			Diese Regel gilt jedoch nicht nur in abstrakten Ersatzschaltbildern, sondern ist auch auf reale, geometrische Formen anwendbar.
			Das rechts dargestellte Beispiel verdeutlicht dies. Fließt ein Strom I in einen Draht hinein, verteilt sich dieser als Stromdichte J im Material. Die Gesamtbilanz über die Fläche ist somit null.
			Mathematisch sagt man, dass das Flächenintegral der Stromdichte über eine geschlossene Fläche immer gleich Null ist.}
			\silence{4}
	}
\end{frame}





\begin{frame}
	\ftx{Example: Parallel connection of resistors}
	
	\subsection{Application case node rule: Parallel connection of resistors}
	
	\s{
	
		A basic application for the node rule is the parallel connection of resistors.
		The current divides at a common node, and a partial current flows through each of the individual resistors.
		After passing through the respective resistors, the partial currents rejoin and continue to flow as a total current.
		This is illustrated in example \ref{expl:knoten}.
		
		\begin{expl}{Parallel connection of resistors}{knoten}
			
				In a parallel circuit, the currents split at the common node.
				How large is the current $I_3$ in relation to the currents $I_0$, $I_1$ and $I_2$?\\
			
			%	\begin{figure}
			\begin{center}
				\input{Tikz/src/2_kirchhoff/Parallelschaltung_Widerstand_Skript}
			\end{center}
			%\end{figure}
			\begin{equation*}
				I_0-I_1-I_2 = 0
			\end{equation*}
			\begin{equation*}
				\rightarrow	I_0 = I_1 + I_2
			\end{equation*}
			\begin{equation*}
				-I_3+I_1+I_2 = 0
			\end{equation*}
			\begin{equation*}
				\rightarrow I_1+I_2 = I_3
			\end{equation*}
			\begin{equation*}
				\rightarrow I_3 = I_0
			\end{equation*}
			
			
			
			
		\end{expl}
		
		If one of the resistors is replaced by an ideally conductive connection (conductance approaches
		infinity, resistance consequently approaches 0), the current flow in the area
		between the two nodes changes (example \ref{expl:knoten2}).. 
		
		
		
		\begin{expl}{Parallel connection with conductor}{knoten2}
			
			A resistor is replaced by a conductive connection.
			How large are the currents $I_1$ and $I_2$?\\
			
			\begin{center}
				\input{Tikz/src/2_kirchhoff/Parallelschaltung_Leiter_Skript}
			\end{center}
			
			\begin{equation*}
				I_0=I_1+I_2
			\end{equation*}
			\begin{equation*}
				\mathrm{mit} \, U_1 = R_1 \cdot I_1
			\end{equation*}
			\begin{equation*}
				U_2 = R_2 \cdot I_2 =0 \stackrel{!}{=} U_1
			\end{equation*}
			\begin{equation*}
				\rightarrow I_1 = 0
			\end{equation*}
			\begin{equation*}
				\rightarrow I_0 = I_2 = I_3
			\end{equation*}
			
			
			
			
			
			
		\end{expl}
	}
	\b{
	
		\begin{columns}
			\column[t]{0.6\textwidth}
			
			
			
			
			%\vspace{-20pt}
			
			In a parallel circuit, the currents split at the common node.
			How large is the current $I_3$?	\\
			
			
			\column[t]{0.4\textwidth}
			
			\input{Tikz/src/2_kirchhoff/Parallelschaltung_Widerstand_Folien}
			
		\end{columns}
		\begin{columns}
			\column[t]{0.6\textwidth}
			\vspace{-50pt}
			
			
			
			
			\pause
			
			
			$I_0-I_1-I_2 = 0$\\
			$I_0 = I_1 + I_2$\\
			$I_1+I_2 = I_3$\\
			$\rightarrow I_3 = I_0$
			
			\pause
			
			\column[t]{0.4\textwidth}
			
			\phantom{.}
			
		\end{columns}
		
		\phantom{.}\\
		\begin{columns}
			\column[t]{0.6\textwidth}
			\vspace{-10pt}
			
			
			A resistor is replaced by a conductive connection.
			How strong is the current? $I_2$?\\
			
			\column[t]{0.4\textwidth}
			
			\input{Tikz/src/2_kirchhoff/Parallelschaltung_Leiter_Folien}
			
		\end{columns}
		\begin{columns}
			\column[t]{0.6\textwidth}
			
			
			\pause
			\vspace{-55pt}
			$I_0 = I_1 + I_2$\\
			
			$U_2 = R_2 \cdot I_2 = 0 \cdot I_2$\\
			$U_1 = R_1 \cdot I_1  \stackrel{!}{=} U_2 =0 $\\
			$\rightarrow I_1 =0$\\
			$\rightarrow I_2 = I_0 = I_3$
			
			\column[t]{0.4\textwidth}
			\phantom{.}\\
			
			
		\end{columns}
		
		\speech{folie12}{1}{Ein klassischer Fall für die Anwendung der Knotenregel ist die Parallelschaltung von Widerständen. Hierbei stellt sich zum Beispiel die Frage wie groß
			der Strom I drei ist.}
			\silence{2}
		\speech{folie12}{2}{Zur Beantwortung dieser Frage müssen die Knotengleichungen aufgestellt werden. Der Gesamtstrom I null teilt sich am Knoten auf,
			ein Teilstrom fließt durch den Widerstand R eins, der andere durch R zwei.  
			Im Knoten unterhalb der Widerstände fließen die Ströme I eins und Zwei in den Knoten hinein, während I drei aus ihm herausfließt. 
			Aus der Anwendung der Knotenregel auf beide Knoten ergibt sich somit I null ist gleich I eins plus I zwei, was wiederum gleich I3 ist. }
			\silence{2}
		\speech{folie12}{3}{Jetzt wollen wir uns anschauen, was passiert, wenn einer der Widerstände durch einen idealen Leiter ersetzt wird.}
		\silence{2}
		\speech{folie12}{4}{Da der Widerstand R eins durch diese leitende Verbindung kurzgeschlossen ist, fließt der gesamte Strom durch die neue Verbindung.
			Somit ist der Strom Ih eins durch den Widerstand null.
			Eingesetzt in die Knotengleichung ergibt sich also, dass
			Ih null gleich Ih zwei gleich Ih drei ist.}
			\silence{10}
		
	}
\end{frame}




\begin{frame}
	\subsection{Mesh rule (Kirchhoff's second law)}
	\s{
	
		Kirchhoff's second law (mesh rule) states that the \textbf{sum of all voltages} in a mesh is \textbf{zero}.
		Analogous to the direction of the currents, the direction of the individual partial voltages must also be taken into account here.
		If the direction arrow of a partial voltage points in the opposite direction to the direction of circulation of the mesh, then
		this partial voltage must be given a negative sign. If the direction of an applied voltage
		is not known, then an arbitrary counting arrow direction can also be assumed here. An oppositely 
		applied voltage is also expressed in calculations with a negative sign. 
		
		\begin{Key point}{Mesh rule}
			
			\begin{equation}
				\sum_{k=1}^{n} U_k = 0
			\end{equation}
			
			
		\end{Key point}
		
		Equivalent to the above definition, it can be stated that the sum 
        of all voltages connected in the same direction at voltage sources corresponds to the voltage that
        drops across the consumers.
		
		\begin{Key point}{Mesh rule 2}
			Total voltage at voltage sources = total voltage at consumers
		\end{Key point}
		
		Applied to the example network shown in Figure \ref{fig:masche}, the following mesh equations result:
		
		\begin{equation*}
			U_1+U_2-U_3-U_4 =0
		\end{equation*}
		respectively
		
		
		
		\begin{equation*}
			U_1+U_2 = U_3+U_4
		\end{equation*}
		%$U_1+U_2 = U_3+U_4$\\
		%$U_1+U_2-U_3-U_4 =0$\\
		
		
		
		
		
		
		
		
		
		
		\begin{figure}[h!]
			\renewcommand{\figurename}{Figure}
			\begin{center}

				\input{Tikz/src/2_kirchhoff/Netzwerk_Spannung_Wiederstaende}
				
			\end{center}
			\caption{Example network consisting of two voltage sources and two resistors}
			\label{fig:masche}
		\end{figure}
		
		The mesh equation also applies if additional currents are fed into the mesh,
        or if individual two-poles are passed through several times during the circulation around a closed mesh.
		
		The general relationship, whereby any integrated voltage along a closed contour equals 0, can
        be described as follows according to Maxwell's equations:
 
		
		\begin{equation}
			\oint_{s} \vec{E} d\vec{s} = 0
		\end{equation}
		
		
		
	}
	
	
	\b{
		\ftx{Mesh rule (Kirchhoff's second law)}
		\begin{columns}
			\column[t]{0.65\textwidth}
			
			\vspace{-70pt}
			
			With closed mesh circulation:
			\vspace{-5pt}
			%\hspace{-40pt}
			\begin{Key point}{}
				
				\begin{center}

					Sum of voltage sources\\
					\quad \quad\ \quad =\\
					Total voltage at consumers
					
				\end{center}
			\end{Key point}
			
			\column[c]{0.35\textwidth}
			
			\vspace{10pt}
			%    \vspace{60pt}
			%    \hspace{30pt}
			
			\input{Tikz/src/2_kirchhoff/Netzwerk_Spannung_Wiederstaende}
			
			\vspace{10pt}
			
			
			$U_1+U_2 = U_3+U_4$\\
			
		\end{columns}
		
		\pause
		\begin{columns}

			\column[t]{0.65\textwidth}
			
			\vspace{-25pt}
			
			
			Or:\\
			With the mesh circulation closed, add up the partial stresses with the correct sign.
			
			\vspace*{-20pt}
			
			\begin{Merksatz}{}
				$\sum_{k=1}^{n} U_k = 0$
			\end{Merksatz}
			
			\column[t]{0.35\textwidth}
			
			
			
			
			
			
			
			$U_1+U_2-U_3-U_4 =0$\\
			
			
			\phantom{.}\\
			
			\pause
			
			General description according to Maxwell's equations:
			
			\vspace{10pt}
			
			$\oint_{s} \vec{E} d\vec{s} = 0$
			
			
		\end{columns}
		\speech{folie13}{1}{Die zweite Kirchhoffsche Regel ist die Maschenregel. Sie besagt
			dass die Summe aller Spannungsquellen in einer geschlossenen Masche gleich der Summe an Spannungen über die Verbraucher ist.
			Jede Spannung, die irgendwo abfällt, muss also von einer Spannungsquelle bereitgestellt werden.
			Im Beispielnetzwerk ergibt sich somit U eins Plus U zwei gleich U drei Plus U vier.}
			\silence{2}
		\speech{folie13}{2}{
			Alternativ kann man bei einem geschlossenem Maschenumlauf die Teilspannungen auch vorzeichenrichtig aufsummieren.
			Also alle Spannungen in Umlaufrichtung der Masche, in unserem Beispiel hier gegen den Uhrzeigersinn, positiv zählen,
			die Spannungen entgegen der Umlaufrichtung hingegen negativ.
			Generell kann die Umlaufrichtung sowie die Richtung der Spannung dabei frei gewählt werden.
			Auch hier gilt wie beim Knotensatz, dass eine in der realität andersherum anliegende Spannung beim Berechnen lediglich ein negatives Vorzeichen erhält.}
		
		\speech{folie13}{3}{Die Maschenregel lässt sich auch physikalisch durch Maxwells Gleichungen ausdrücken.
			Die aufintegrierte elektrische Feldstärke entlang eines geschlossenen Weges ergibt null, wobei die genaue Herleitung dieses Zusammenhanges über den Umfang dieser Einführung hinaus geht.}
			\silence{4}
			}
	
\end{frame}

%\begin{frame}
%	\fta{Testbild}
%	\begin{circuitikz}
%		% Zeichne die Spannungsquellen in Reihe
%		\draw (0,0) to[battery1, l_=$V_1$] (0,3)
%				   to[battery1, l_=$V_2$] (0,6);
%	
%		% Verbinde die Spannungsquellen mit den Widerständen durch Leitungen
%		\draw (0,6) -- (4,6);
%		\draw (0,0) -- (4,0);
%	
%		% Zeichne die Widerstände in Reihe
%		\draw (4,0) to[R, l_=$R_1$] (4,3)
%				   to[R, l_=$R_2$] (4,6);
%	
%		% Zeichne den kreisförmigen Pfeil in der Mitte
%		\draw[->, thick] (2,5) arc[start angle=90, end angle=-270, radius=1]
%		\node at (2,2) {Masche};
%	\end{circuitikz}
%\end{frame}	





\begin{frame}

	\subsection{Application case mesh rule: Series connection of resistors}
	
	\s{
	
		While the nodes in electrical networks are primarily important for current, the branches and thus also the meshes are primarily of interest for
		calculating voltages. 
		In a series connection of resistors in a branch, all partial voltages add up to a total voltage with the correct sign.
        An application case for determining a partial voltage is shown in Example \ref{expl:reihe}.
		
		
		
		\begin{expl}{Series connections in networks}{reihe}
			
			
			In a series connection, the voltages add up to a
			total voltage.
			How large is the voltage $U_0$?
			
			\begin{center}

				\input{Tikz/src/2_kirchhoff/Netzwerk_Reihenschaltung}
				
			\end{center}
			
			Set up the mesh counterclockwise:
			
			\begin{equation*}
				U_2-U_0-U_3 =0
			\end{equation*}
			
			\begin{equation*}
				\rightarrow U_0 = U_2- U_3
			\end{equation*}
			
			
			
			
		\end{expl}
		
		
        The mesh rule can also be applied if the mesh to which it is applied has an interruption point (see example \ref{expl:reihe2}). 

		\begin{expl}{Series connections with interruption point}{reihe2}
			
			
			A connection is interrupted. How large is the voltage $U_1$ at the point of interruption?
			
			\begin{center}

				\input{Tikz/src/2_kirchhoff/Netzwerk_Reihenschaltung_Unterbrechung}
				
			\end{center}
			
			Set up the mesh counterclockwise:
			
			%\begin{equation*}
			%U_2+U_1-U_0-U_3 = 0
			%\end{equation*}
			
			\begin{equation*}
				U_2+U_1-U_0-U_3 = 0
			\end{equation*}
			
			Determine the voltage $U_2$ using Ohm's law:
			
			
			\begin{equation*}
				U_2=I_2 \cdot R_2
			\end{equation*}
			
			\begin{equation*}
				\rightarrow I_2 = 0
			\end{equation*}
			
			\begin{equation*}
				\rightarrow U_2 = 0
			\end{equation*}
			
			Insert into the initial equation and solve for $U_1$: 	
			
			\begin{equation*}
				U_1 = U_0 + U_3
			\end{equation*}

			
		\end{expl}
		
	}
	
	\newpage
	
	\b{
		
		\ftx{Example: Series connection of resistors}
		\begin{columns}
			\column[t]{0.6\textwidth}
			
			\vspace{-50pt}
			In a series connection, the voltages add up to a
			total voltage.\\
			How high is the voltage? $U_0$?
			
			\column[c]{0.4\textwidth}
			
			\input{Tikz/src/2_kirchhoff/Netzwerk_Reihenschaltung}

		\end{columns}
		
		\begin{columns}
			\column[t]{0.6\textwidth}
			
			\pause
			
			\vspace{-30pt}
			$U_2-U_0-U_3 =0$\\
			$\rightarrow U_0 = U_2- U_3$
			
			\column[t]{0.4\textwidth}
			
			\phantom{.}\\
			
		\end{columns}
		
		\begin{columns}
			\column[t]{0.6\textwidth}
			\pause
			
			A connection is interrupted. How large is the voltage $U_1$ at
			the point of interruption?
			
			\column[t]{0.4\textwidth}
			
			\input{Tikz/src/2_kirchhoff/Netzwerk_Reihenschaltung_Unterbrechung}
			
		\end{columns}
		
		\pause 
		
		\begin{columns}
			\column[t]{0.6\textwidth}
			
			\vspace{-55pt}
			
			$U_2+U_1-U_0-U_3 = 0$\\
			$U_2=I_2 \cdot R_2$\\
			$\rightarrow I_2 = 0$\\
			$\rightarrow U_2 = 0$\\
			$U_1 = U_0 + U_3$
			
			\column[t]{0.4\textwidth}
			
		\end{columns}
		
		\speech{folie14}{1}{Ein typischer Anwendungsfall der Maschenregel ist die Reihenschaltung von Widerständen. Hier kann beispielsweise eine Teilspannung gesucht sein.
			In diesem Beispiel wird die Größe der Spannung U null bestimmt.}
			\silence{2}
		\speech{folie14}{2}{Angenommen wir stellen die Masche entgegen dem Uhrzeigersinn auf, so folgt, U zwei minus U null minus U drei gleich null.
			Das bedeutet, U null ist gleich U zwei minus U drei.}
			\silence{2}
		\speech{folie14}{3}{Nun wird eine Verbindung unterbrochen. Gesucht ist infolgedessen die Spannung U1 an der Unterbrechungsstelle.}
		\silence{2}
		\speech{folie14}{4}{Auch hier bilden wir die Masche entgegen dem Uhrzeigersinn. So ergibt sich, U zwei Plus U eins minus U null minus U drei ist gleich 0.}
		\silence{2}
		\speech{folie14}{4}{Da durch die geöffnete Leitung kein Strom mehr fließt, gilt an dieser Stelle,
			I gleich 0. 
			Dem Ohmschen Gesetz U gleich R mal I zu Folge ist U zwei gleich 0, weil I zwei gleich 0 ist. Durch Einsetzen in die Ausgleichsgleichung erhalten wir das Ergebnis, 
			dass U eins gleich U null Plus U drei ist.
			Als Fazit können wir festhalten:
			Die Maschenregel funktioniert also auch bei Stromunterbrechungen beziehungsweise offenen Leitungen.}
			\silence{4}
	}
	
	
\end{frame}
